%% -*- coding:utf-8 -*-
%%%%%%%%%%%%%%%%%%%%%%%%%%%%%%%%%%%%%%%%%%%%%%%%%%%%%%%%%
%%   $RCSfile: 5-hpsg-semantik.tex,v $
%%  $Revision: 1.22 $
%%      $Date: 2008/09/30 09:14:41 $
%%     Author: Stefan Mueller (CL Uni-Bremen)
%%    Purpose: 
%%   Language: LaTeX
%%%%%%%%%%%%%%%%%%%%%%%%%%%%%%%%%%%%%%%%%%%%%%%%%%%%%%%%%

\chapter{Semantik}
\label{Kapitel-sem}\label{Kapitel-Semantik}


Innerhalb\is{Semantik|(} der HPSG"=Forschung gibt es drei verschiedene Ansätze
zur Beschreibung der Bedeutung sprachlicher Ausdrücke. In den ersten
HPSG"=Arbeiten \citep{ps,ps2,Mueller99a} und auch in einigen Umsetzungen 
innerhalb der Computerlinguistik wurde die Situationssemantik
\citep*{BP83a,CMP90,Devlin92}\nocite{BP87a} verwendet.
Eine Publikation neueren Datums, die ebenfalls Situationssemantik verwendet,
ist \citew{GSag2000a-u}. Arbeiten neueren Datums nutzen sogenannte
Unterspezifikationssemantik. Dazu zählen die \emph{Lexical Resource Semantics} (LRS,
\citealt{RS2004a-u}) und die \emph{Minimal Recursion Semantics}\is{Minimal Recursion Semantics@\emph{Minimal Recursion Semantics} (MRS)}
(MRS) \citep*{CFPS2005a}. Während die LRS hauptsächlich in Frankfurt/Main verwendet wird, ist MRS
weiter verbreitet. MRS wird in theoretische Arbeiten (\zb \citew{Kiss2001a}) benutzt, aber auch in
Computerimplementationen und Übersetzungssystemen \parencites{BFO2002a-u}{MuellerCoreGram}. MRS ist für computerlinguistische Systeme gut geeignet, da
Skopusbeziehungen\is{Skopus} unterspezifiziert dargestellt werden können, was eine effizientere
Verarbeitung der Grammatiken durch Computer ermöglicht.
Das folgende Kapitel führt in MRS ein.

Dieses Buch ist ein Syntax"=Buch. Syntax ist ein komplexes System und die Beschäftigung damit
faszinierend. Allerdings ist Syntax allein im wahrsten Sinne des Wortes sinnlos. Die (Morpho-)Syntax wird nur gebraucht, um eine
Verbindung zwischen der lautlichen Form einer Äußerung und ihrer Bedeutung herzustellen: Eine
Höhrer*in ist nur in der Lage herauszufinden, was eine Sprecher*in mit ihrer Äußerung meint, wenn
sie in der Lage ist, die Laute in Wörter und Wortgruppen zu strukturieren und einzelnen Wörtern,
größeren Einheiten und schließlich dem gesamten komplexen Gebilde
eine Bedeutung zuzuordnen, \dash aus der Bedeutung einzelner Wörter, die Bedeutung der gesamten
Äußerung zu ermitteln. Wenn es die Aufgabe der Syntax ist, entsprechende Strukturen bereitzustellen,
muss man als Grammatiker*in auch sagen können, wie Bedeutung repräsentiert wird und wie die
Bedeutungskomposition mittels syntaktischer Strukturen funktioniert. In jedem Grammatikframework
muss es dazu Überlegungen geben. Verglichen mit anderen Theorien wird bei vielen HPSG"=Arbeiten über Syntax und Semantik eines
bestimmten Phänomens gesprochen.

\section{Relationen und ihre Repräsentation in Merkmalbeschreibungen}





\section{Linking}


\section{Komposition}

\section{Quantorenskopus}




%\section*{Kontrollfragen}

\questions{
\begin{enumerate}
\item Was versteht man unter Linking?
\item Wie wird in der HPSG eine Verbindung zwischen Form und Bedeutung eines (komplexen)
      sprachlichen Zeichens hergestellt?
\end{enumerate}
}

%\section*{Übungsaufgaben}

\exercises{
\begin{enumerate}
\item Wie kann man den semantischen Beitrag von \emph{lacht} repräsentieren? Geben Sie den
  Lexikoneintrag für \emph{lacht} an und berücksichtigen Sie das Linking.
\item Geben Sie die vollständige Merkmalbeschreibung für die Wortfolge in (\mex{1}), wie sie in
  \emph{dass er lacht} vorkommt,  an: 
      \ea
      er lacht
      \z
      Dabei soll sowohl die syntaktische Struktur als auch der Bedeutungsbeitrag berücksichtigt werden.

\item Laden Sie die zu diesem Kapitel gehörende Grammatik von der Grammix"=CD
(siehe Übung~\ref{uebung-grammix-kapitel4} auf Seite~\pageref{uebung-grammix-kapitel4}).
Im Fenster, in dem die Grammatik geladen wird, erscheint zum Schluß eine Liste von Beispielen.
Geben Sie diese Beispiele nach dem Prompt ein und wiederholen Sie die in diesem Kapitel besprochenen
Aspekte.
\end{enumerate}
}




\furtherreading{
\citet{KR2024a} geben einen Überblick über Semantik im Rahmen der HPSG. \citet{DKW2024a}
beschäftigen sich mit Argumentstruktur und Linking.

Der Artikel von
\citet{CFPS2005a} stellt MRS vor. Ich kann mich erinnern, dass ich in den 90ern diesen Artikel nicht
verstehen konnte, obwohl ich an einer Grammatik gearbeitet habe, die eine MRS enthielt. Ich habe
damals gemeinsam mit Walter Kasper an der deutschen \Verbmobil"=Grammatik gearbeitet. Das Buch von
\citet{BB2005a} hat mir sehr dabei geholfen, die Grundidee von Semantikframeworks mit
Unterspezifikation zu verstehen. Das Buch kommt mit einer Prolog"=Implementation der entsprechenden
Grammatikfragmente, so dass man Dinge selbst ausprobieren kann. Ich hoffe, dass das hier vorliegende
Kapitel verständlich und einfach genug ist. Wenn das nicht der Fall ist, können vielleicht
\citeauthor{BB2005a} helfen oder der Originalartikel von \citeauthor{CFPS2005a}.

}